% Section 18.11, from lectures/L18/appendix/b_exercises.md.

\section{Exercises}
\label{c18:sec:exercises}

\begin{aside}
Write your own \file{register_bank.vhd} from \secref{c18:sec:iface} to \ref{c18:sec:write}; the
testbench handed out, run per appendix~\ref{app:sim}, is the check that it behaves correctly.
Exercise~\ref{c18:ex:waveform} and \ref{c18:ex:slave} are about \file{spi_slave.vhd}, which is handed
out rather than written \textemdash{} that one you read, and \chapref{c20:ch} expects you to know it.

The authoritative source for every register number below is appendix~\ref{app:protocol}. Where an
exercise and a specification disagree, the specification wins, and the exercise is the bug.
\end{aside}

\exercise{\code{register_bank}}{VHDL}
\label{c18:ex:bank}
Write your own \file{register_bank.vhd} from the chapter. Verify it:

\begin{shell}
cd bridge
ghdl -a --std=93 ../controller/can_def.vhd register_bank.vhd register_bank_tb.vhd
ghdl -e --std=93 register_bank_tb
ghdl -r --std=93 register_bank_tb --assert-level=error
\end{shell}

If a check fails, read the assertion message; each of the eleven cases says what a failure in that
particular case means.

\exercise{The three STATUS bits}{Reasoning}
\label{c18:ex:status}
\xpart{a} Out of reset, STATUS reads \code{0x00000001}. Say which bit that is and why the reset value
is \code{'1'} and not \code{'0'}. Then describe what a driver's \code{send()} would do, forever, if it
were \code{'0'}.

\xpart{b} Bit 2 latches a \textbf{rising edge} of the controller's \code{error}, not its level.
Suppose it latched the level instead. Write down the sequence of driver operations \textemdash{} poll,
clear, poll \textemdash{} that would then behave differently, and say what the driver's author would
conclude was broken.

\xpart{c} Bit 1 is set by \code{rx_valid} and cleared by a write to \code{RX_ACK}. Bit 0's main set
condition is \code{tx_done}, and it is cleared by a write to \code{TX_SEND}. One of those pairs has
its set and clear conditions in the ``natural'' order, and the other is reversed. Say which is which,
and why the reversed one is right all the same.

\xpart{d} \Secref{c18:sec:status} gives an ordering rule: controller events are applied \textbf{after}
the bridge's writes in the clocked process. Construct the one clock edge where that rule changes the
outcome, name both events, and say which of the two possible resolutions can lock the system up and
which cannot.

\xpart{e} Bit 0 has \emph{three} set conditions, and one of them is a rising edge of \code{error}
while bit 0 is already low. Explain why that condition needs to know nothing about whether the
controller was transmitting or receiving \textemdash{} what is it about a node's behaviour that makes
``error rose while we were busy'' mean ``a transmission has just ended''?

\xpart{f} Remove that condition from your bank and rerun the testbench. Which case fails, and what
consequence does its message say it would have had? Then explain why this is defence in depth rather
than duplicated work, given that \chapref{c17:ch} already has the controller pulse \code{tx_done} at
an aborted transmission. (Who writes \code{can_controller}? Who wrote the testbench that would catch
it?)

\xpart{g} \code{TX_ABORT} is the third path, and unlike the other two it does not correspond to
anything the hardware did \textemdash{} it is the driver going around the bank. Say what it does
\emph{not} do: name at least two things you might expect a register called \code{TX_ABORT} to reach,
and explain why it reaches neither. (The clue is \code{can_controller}'s port list, settled in
\chapref{c11:ch}.)

\exercise{\code{TX_SEND} and the pulse discipline}{Simulation}
\label{c18:ex:txsend}
\xpart{a} A write of \code{0x1} to \code{TX_SEND} raises \code{tx_req} for exactly one cycle. Say
exactly which cycle, relative to the clock edge that carried \code{reg_write}, and which two things
\code{register_bank_tb} checks about it.

\xpart{b} Now break it deliberately. Change your \code{TX_SEND} branch so that \code{tx_req} becomes a
\textbf{level}: set it at an accepted write and clear it only when \code{tx_done} arrives. Rerun
\code{register_bank_tb}. Which check fails first, and what does its message say?

\xpart{c} Your broken version from \textbf{b)} would still \emph{look} correct on an oscilloscope
watching a single frame. Describe, using \secref{c16:sec:roles}'s rule for when
\code{can_controller} acts on \code{tx_req}, what it does at the second frame and every frame after
it. Then say why this is a bug simulation catches easily and a test on the bench does not.

\xpart{d} A write to \code{TX_SEND} while STATUS bit 0 is low is discarded entirely \textemdash{} no
pulse, no state change. Another design could queue it instead and fire when the controller reports
ready. Give a concrete argument for the discard rule, given that the driver can poll STATUS. Put the
module back as it was.

\exercise{Reading an SPI transaction off the waveform}{Reasoning}
\label{c18:ex:waveform}
This chapter is called ``from the waveform'' for a reason: \chapref{c20:ch} may well hand you one and
ask what it means. Mode 0 throughout \textemdash{} SCK idles low, MOSI is sampled on the
\textbf{rising} edge, MISO changes on the \textbf{falling} edge, MSB first.

\xpart{a} Draw by hand the whole five-byte transaction writing \code{0x00000001} to \code{TX_SEND}.
Work out the command byte from the format in appendix~\ref{app:protocol} rather than copying it: bit 7
is the write bit, bits 6--4 are reserved, bits 3--0 are the register index, and \code{TX_SEND} is
index 5. Label \code{SS}, \code{SCK} and \code{MOSI}, and mark the 40 rising edges in groups of eight.

\xpart{b} Mark on the same drawing where the master must not release \code{SS}, and what the slave has
to do if it does anyway. Restate the rule from the protocol specification's section on aborts in your
own words.

\xpart{c} A read of \code{STATUS} is the mirror image: command byte \code{0x00}, four dummy bytes out,
four bytes of register value back on MISO. The protocol specification says the slave latches the
register value \textbf{once}, at the end of the command byte. Explain what a driver could observe if
the slave instead reread the register for each of the four bytes, and why \code{STATUS} in particular
is the register where that would hurt.

\xpart{d} Count the SCK edges in a whole transaction and, at the maximum of 1~MHz the protocol
specification allows, work out how long a transaction takes. Compare that with the 20~ns pulse the
register bank exists to translate. That ratio is the whole argument for this layer, in a single
number.

\exercise{Reading \file{spi_slave.vhd}}{Reasoning}
\label{c18:ex:slave}
You do not write this module, but you have to be able to read it.

\xpart{a} \code{spi_slave} does not use \code{SCK} as a clock. It synchronises it into the 50 MHz
domain and detects its edges instead. Give two separate reasons, one about metastability
(\chapref{c4:ch}) and one about what a design with two unrelated clocks costs in an FPGA.

\xpart{b} The record \code{sync_t} has three fields, but only two of them are synchronisation stages.
Say what the third is for, and why \code{mosi} needs two fields where \code{sclk} and \code{ss} need
three.

\xpart{c} In the reset branch, \code{sclk_s} is cleared to all zeros and \code{ss_s} to all
\textbf{ones}. Explain why they differ, and then trace what \code{ss_active} would read during the
first two clock edges after reset if \code{ss_s} were cleared to zeros with no master connected. This
is the same rule Exercise~\ref{c4:ex:buttonsync} was built around; say which check in that exercise's
testbench corresponds to it.

\xpart{d} The SCK edge handling sits inside the \code{else} arm of a test on \code{ss_s.s2}, so a
deselected slave ignores the SPI lines entirely. Construct the failure that would follow if it did
not: start from a stray SCK edge while \code{SS} is high, and follow it through to what the
\emph{next} register write lands on.

\xpart{e} The \code{ss_falling} branch is an \code{elsif} before the SCK branches rather than a
separate \code{if} beside them. Say at which single clock edge that order matters, and what would be
wrong with the received bytes if the two branches ran in the reverse order.

\exercise{The capture, and the gap it does not close}{Reasoning}
\label{c18:ex:capture}
\xpart{a} When \code{rx_valid} pulses, the bank captures \code{rx_id}, \code{rx_dlc} and
\code{rx_data} into the \code{RX_*} registers instead of wiring the controller's ports straight into
the read multiplexer. Explain what a driver reading four separate SPI transactions would see without
that capture.

\xpart{b} A write to \code{RX_ACK} clears STATUS bit 1 but leaves the \code{RX_*} registers holding
their contents. Say why that is harmless, and quote the sentence in the register map that makes it so.

\xpart{c} A second frame arrives before the first has been acknowledged. Say exactly what the driver
observes, and why this is a documented gap inherited from \code{can_controller} (\chapref{c11:ch})
rather than a shortcoming of the register bank.

\xpart{d} Design, on paper alone, the smallest change to the \emph{register map} that would let a
driver detect that it had missed a frame. You are not adding a receive queue; you are adding the
ability to know that one was needed. Say what it costs in hardware, and in which existing register it
could live without breaking the contract the parallel class is already writing to.
